\documentclass[12pt, fleqn]{article}
\usepackage[portuguese]{babel}
\usepackage{newpxtext, newpxmath}
\usepackage{amsmath}
\usepackage{titling}
\usepackage{enumitem}
\usepackage{geometry}
\usepackage{graphicx}
\usepackage{float}
\usepackage{listings}
\usepackage{xcolor}
\usepackage{booktabs} % Pacote para tabelas profissionais

\geometry{
    paper=a4paper,
    top=64pt,
    bottom=64pt,
    left=112pt,
    right=112pt
}

\lstset{
    basicstyle=\ttfamily\small,
    backgroundcolor=\color{gray!5},
    rulecolor=\color{gray!20},
    frame=single,
    framesep=16pt,
    xleftmargin=16pt,
    xrightmargin=16pt,
    captionpos=b,
    breaklines=true
}

\renewcommand{\arraystretch}{1.2}

\setlength{\tabcolsep}{8pt}

\NewDocumentCommand{\course}{m}{\gdef\thecourse{#1}}
\NewDocumentCommand{\authors}{m}{\gdef\theauthors{#1}}
\NewDocumentCommand{\professor}{m}{\gdef\theprofessor{#1}}
\NewDocumentCommand{\semester}{m}{\gdef\thesemester{#1}}

\renewcommand{\maketitle}{
    \begin{titlepage}
        \begin{center}
            {\scshape Universidade Federal do Rio Grande do Sul\\Instituto de Informática} \\
            \vspace*{\fill}
            {\large \thecourse} \\[16pt]
            {\LARGE \thetitle} \\[36pt]
            \theauthors \\
            \vspace*{\fill}
            \theprofessor \\[2pt]
            \thesemester
        \end{center}
    \end{titlepage}
}

% --- INFORMAÇÕES DO DOCUMENTO ATUALIZADAS ---
\course{Circuitos Digitais – INF01058}
\title{LAB10 – Unidade de controle do Neander}
\authors{
    Ana Cláudia Rodrigues — 343123 \\[2pt]
    Bruno Samuel Ardenghi Gonçalves — 550452
}
\professor{Prof. Mateus Grellert}
\semester{2025/1}

\begin{document}

\maketitle

% --- SEÇÃO DE INTRODUÇÃO ATUALIZADA ---
\section{Introdução}

Neste laboratório, foi projetada a Unidade de Controle (UC) do processador Neander, um componente essencial do projeto final da disciplina. A UC é responsável por gerar todos os sinais que governam o datapath do processador, orquestrando a execução de cada instrução.

O controle foi implementado de forma "hardcoded", utilizando um conjunto de expressões lógicas booleanas para cada sinal de controle, em vez de uma máquina de estados finitos explícita. Para determinar o estado de cada sinal, a lógica combinacional da UC utiliza como entradas os bits da instrução atual, fornecidos por um decodificador de opcode, e o passo de execução atual (t0-t7), fornecido por um módulo temporizador.

O projeto foi implementado e simulado no software Quartus, integrando os módulos UC, Decodificador e Temporizador para validação funcional.

\subsection{Entradas}

As entradas principais para o sistema de controle como um todo são:
\begin{itemize}
	\item \texttt{CLK}: Clock principal do sistema.
	\item \texttt{RESET}: Sinal para reiniciar o estado do temporizador e do sistema.
	\item \texttt{OPCODE[3..0]}: Código da instrução atual, vindo do Registrador de Instrução (RI), que serve de entrada para o Decodificador.
	\item \texttt{N, Z}: Flags de estado (Negativo e Zero) da UAL, utilizadas para a execução de saltos condicionais.
	\item \texttt{HLT}: Sinal que paralisa a contagem do temporizador, efetivamente parando o processador.
\end{itemize}

\subsection{Saídas}

As saídas são os 17 sinais de controle que governam os componentes do datapath do Neander. Os principais são:
\begin{itemize}
	\item \texttt{cargaPC, incrementaPC}: Sinais de controle do Contador de Programa (PC).
	\item \texttt{cargaAC, cargaNZ}: Sinais de controle do Acumulador (AC) e dos registradores de flag.
	\item \texttt{cargaRI, cargaREM, cargaRDM}: Sinais de habilitação para os registradores de Instrução, Endereço de Memória e Dados de Memória.
	\item \texttt{selUAL[2..0]}: Sinal de 3 bits que seleciona a operação a ser realizada pela UAL.
	\item \texttt{Read, Write}: Sinais para controlar as operações de leitura e escrita na memória.
\end{itemize}

\subsection{Módulos Principais}

\begin{itemize}
	\item \texttt{Unidade de Controle (UC)}: O núcleo do projeto. Um circuito puramente combinacional que implementa as expressões lógicas para gerar os sinais de controle.
	\item \texttt{Decodificador}: Recebe o \texttt{OPCODE[3..0]} e gera um sinal de saída para cada uma das 11 instruções do Neander.
	\item \texttt{Temporizador}: Um contador de 3 bits que gera 8 sinais de tempo distintos (\texttt{t0} a \texttt{t7}), sequenciando os passos de execução de cada instrução.
\end{itemize}

% --- SEÇÃO DE RESULTADOS COM TABELAS ATUALIZADAS ---
\section{Resultados}

\begin{table}[H]
    \centering
    \begin{tabular}{|l | c|}
        \hline
        \textbf{Total de elementos lógicos} & 40 \\
        \hline
        \textbf{Total de pinos} & 24 \\
        \hline
        \textbf{Área total} & 64 \\
        \hline
        \textbf{Total de registradores} & 3 \\
        \hline
    \end{tabular}
    \caption{Consumo de recursos da FPGA em área}
\end{table}

\begin{table}[H]
    \centering
    \begin{tabular}{|l | c|}
        \hline
        \textbf{Frequencia máxima} & 588.24 MHz \\
        \hline
        \textbf{Atraso crítico} & 1.662  \\
        \hline
    \end{tabular}
    \caption{Frequência máxima de operação e atraso crítico (\textit{slow}, 85°C)}
\end{table}

OBS: o atraso crítico é dado 
\begin{itemize}
	\item \texttt{de}: TIMER:TIMER|UD\_COUNTER\_3BIT:inst|inst35
	\item \texttt{para}: TIMER:TIMER|UD\_COUNTER\_3BIT:inst|UD\_COUNTER\_1BIT:inst9|inst7 
\end{itemize}

% --- SEÇÕES GENÉRICAS MANTIDAS ---
\section{Capturas de tela}

\begin{figure}[H]
	\centering
	\includegraphics[width=\textwidth]{netlist.png}
	\caption{Captura de tela do netlist (RTL Viewer)}
\end{figure}

\begin{figure}[H]
	\centering
	\includegraphics[width=\textwidth]{waveform.png}
	\caption{Captura de tela da simulação em forma de onda}
\end{figure}

% --- SEÇÃO DE SDC ATUALIZADA COM MELHORES PRÁTICAS ---
\section{Arquivo de restrições}

\begin{lstlisting}[caption=constraints.sdc, label=lst:sdc]
create_clock -name clk -period 20 [get_ports CLK]

set_input_delay -clock clk -max 0.2 [all_inputs]
set_input_delay -clock clk -min 0.01 [all_inputs]
set_output_delay -clock clk -max 0.2 [all_outputs]
set_output_delay -clock clk -min 0.01 [all_outputs]
\end{lstlisting}

\end{document}